% Generated draft: requires mathematical and source review.
\documentclass[11pt]{article}
\usepackage[margin=25mm]{geometry}
\usepackage{fontspec}
\setmainfont{Latin Modern Roman}
\usepackage{amsmath,amssymb,mathtools,unicode-math}
\setmathfont{Latin Modern Math}
\usepackage{enumitem,needspace,xurl,hyperref}
\setlength{\parindent}{0pt}
\setlength{\parskip}{5pt}
\setlength{\emergencystretch}{3em}
\newcommand{\N}{\mathbb{N}}
\newcommand{\Z}{\mathbb{Z}}
\newcommand{\Q}{\mathbb{Q}}
\newcommand{\R}{\mathbb{R}}
\newcommand{\C}{\mathbb{C}}
\newcommand{\F}{\mathbb{F}}
\newcommand{\A}{\mathbb{A}}
\newcommand{\PP}{\mathbb{P}}
\newcommand{\E}{\mathbb{E}}
\newcommand{\eps}{\varepsilon}
\newcommand{\dd}{\,\mathrm{d}}
\begin{document}
\section*{UnsolvedMath ID: 30005019}
\textbf{Source problem code:} catalogue entry
\subsection{Definitions and mathematical statement}
The supplied question concerns adjoint curves associated with polygonal orders on a fixed set of points in \(\PP^2\). The base field of \(\PP^2\), admissible point configurations, and the meaning of an adjoint curve are not specified in the record.

The combinatorial indexing can be made explicit as follows. For an integer \(n\geq 3\) and distinct points \(P=\{p_1,\ldots,p_n\}\subset\PP^2\), let \(\mathcal O(P)\) denote the set of lists of all points of \(P\), modulo cyclic rotation and reversal. Then
\[
\lvert\mathcal O(P)\rvert=\frac{(n-1)!}{2}.
\]
For \(\omega=[(p_{i_1},\ldots,p_{i_n})]\in\mathcal O(P)\), the corresponding polygon has successive vertices \(p_{i_1},\ldots,p_{i_n}\), with \(p_{i_{n+1}}=p_{i_1}\); its supporting edge lines are \(\overline{p_{i_j}p_{i_{j+1}}}\), for \(1\leq j\leq n\). Whenever this polygon has a uniquely defined adjoint curve in the intended sense, denote that curve by \(A_\omega\).

\textbf{Source question.} For the same fixed point set \(P\), describe the relations among the adjoint curves in the indexed family
\[
\bigl(A_\omega\bigr)_{\omega\in\mathcal O(P)}.
\]
The number \((n-1)!/2\) counts polygonal orders; the record does not assert that their adjoint curves are pairwise distinct.

\textbf{Unresolved qualifications.} The restrictions \(n\geq3\) and distinctness make the displayed indexing and count meaningful; the supplied statement does not explicitly give a range for \(n\) or configuration hypotheses. It provides neither a definition of the adjoint construction nor hypotheses ensuring its existence, uniqueness, and invariance under reversal. In particular, no genericity, noncollinearity, or treatment of degenerate polygons is supplied. The kind of relation sought among these curves is also unspecified. Consequently, the displayed question is a formal description of the intended family, conditional on clarification of these geometric prerequisites, rather than a fully specified conjecture.
\subsection{Short English statement}
Arrange the same \(n\) projective points in every cyclic order, identifying orders that differ by reversal, and ask how the associated polygon adjoint curves are related. There are \((n-1)!/2\) such orders for \(n\geq3\) distinct points, but their curves may coincide. A precise geometric formulation still requires a base field, a definition of the adjoint curve, admissibility hypotheses, and clarification of the relations being sought.
\subsection{Sources}
\begin{itemize}
\item Local UnsolvedMath registry, supplied source record for permanent ID \(30005019\), source problem code \texttt{catalogue entry}, title \emph{Relations among Adjoint Curves from Cyclic Point Orders}. The record belongs to the Oberwolfach Reports Open Problems collection. No proposer is identified; the supplied proposed year is \(2022\).
\item Supplied report citation: \emph{Mini-Workshop: Interpolation, Approximation, and Algebra (2023)}, Oberwolfach Reports, pp.~383--428, DOI \texttt{10.4171/owr/2022/7}, \url{https://doi.org/10.4171/owr/2022/7}. This is also the supplied literature-source link labelled ``OWR adjoints problem session.'' The raw excerpt includes a page marker for p.~425. The citation year \(2023\) and the record's \(2022\) proposed year are retained without reconciliation.
\item Supplied literature source: Agostini--Plaumann--Sinn--Wesner, ``heptagon adjoints,'' \url{https://arxiv.org/abs/2408.15759}. The supplied assessment says that the \(2024\) heptagon result resolves the generic degree of one adjoint map but does not resolve the general relations among curves obtained from different cyclic orders on the same point set.
\item Literature assessment date: \texttt{2026-08-22}. The supplied status is ``open.'' This is supplied dated curation, described in the record as literature triage derived from OWR curation; neither its time-sensitive status nor the cited mathematical claims have been independently verified here.
\item The supplied statement status is \texttt{corrected\_verified}, with a supplied claim that the clean statement was checked against the official report and preserves its hypotheses, notation, and target. This verification claim is recorded as source metadata, not as an independent verification by this draft.
\end{itemize}
\end{document}
